\documentclass[11pt,a4paper]{article}

% ── Packages ──────────────────────────────────────────────────────────
\usepackage[utf8]{inputenc}
\usepackage[T1]{fontenc}
\usepackage{amsmath,amssymb,amsthm,mathtools}
%\usepackage{thmtools}
\usepackage{geometry}
\geometry{margin=1in}
\usepackage{hyperref}
\hypersetup{colorlinks=true,linkcolor=blue!60!black,citecolor=blue!60!black,urlcolor=blue!60!black}
\usepackage{enumitem}
\usepackage{booktabs}
\usepackage{algorithm}
\usepackage{algorithmic}
\usepackage{natbib}
\bibliographystyle{plainnat}
\usepackage{xcolor}
\usepackage{tikz}
\usetikzlibrary{arrows.meta,positioning,decorations.pathreplacing,calc}
\usepackage{caption}
\usepackage{subcaption}
%\usepackage{microtype}

% ── Theorem environments ──────────────────────────────────────────────
\theoremstyle{plain}
\newtheorem{theorem}{Theorem}[section]
\newtheorem{proposition}[theorem]{Proposition}
\newtheorem{lemma}[theorem]{Lemma}
\newtheorem{corollary}[theorem]{Corollary}
\theoremstyle{definition}
\newtheorem{definition}[theorem]{Definition}
\newtheorem{assumption}[theorem]{Assumption}
\newtheorem{example}[theorem]{Example}
\theoremstyle{remark}
\newtheorem{remark}[theorem]{Remark}

% ── Macros ────────────────────────────────────────────────────────────
\newcommand{\R}{\mathbb{R}}
\newcommand{\E}{\mathbb{E}}
\newcommand{\Prob}{\mathbb{P}}
\newcommand{\cO}{\mathcal{O}}
\newcommand{\cA}{\mathcal{A}}
\newcommand{\cX}{\mathcal{X}}
\newcommand{\cY}{\mathcal{Y}}
\newcommand{\cL}{\mathcal{L}}
\newcommand{\cN}{\mathcal{N}}
\newcommand{\cS}{\mathcal{S}}
\newcommand{\cF}{\mathcal{F}}
\newcommand{\cM}{\mathcal{M}}
\newcommand{\cD}{\mathcal{D}}
\newcommand{\cP}{\mathcal{P}}
\newcommand{\deff}{d_{\mathrm{eff}}}
\newcommand{\tr}{\mathrm{tr}}
\newcommand{\diag}{\mathrm{diag}}
\newcommand{\sign}{\mathrm{sign}}
\newcommand{\supp}{\mathrm{supp}}
\newcommand{\poly}{\mathrm{poly}}
\newcommand{\rank}{\mathrm{rank}}
\newcommand{\ip}[2]{\langle #1, #2 \rangle}
\newcommand{\norm}[1]{\| #1 \|}
\newcommand{\abs}[1]{| #1 |}
\DeclareMathOperator*{\argmin}{arg\,min}
\DeclareMathOperator*{\argmax}{arg\,max}
\DeclareMathOperator{\prox}{prox}

% ── Title ─────────────────────────────────────────────────────────────
\title{%
  \textbf{The Blessing of Dimensionality for Learned Representations:}\\[4pt]
  \large A Universal Embedding Theorem Unifying Concentration of Measure,\\
  Superposition, and Causal Feature Recovery
}

\author{
  Anonymous Authors\thanks{%
    Affiliations withheld for review.  
    Correspondence: \texttt{[redacted]}.
  }
}

\date{\today \quad {\small (Working Draft --- Not for Distribution)}}

\begin{document}
\maketitle

% ══════════════════════════════════════════════════════════════════════
% ABSTRACT
% ══════════════════════════════════════════════════════════════════════
\begin{abstract}
The empirical success of large-scale neural networks contradicts the classical
curse of dimensionality: increasing embedding dimensions, parameter counts, and
data simultaneously \emph{improves} rather than degrades performance.  We provide
a unified mathematical explanation rooted in the \emph{blessing of
dimensionality}.  We introduce the \textbf{Universal Embedding Theorem}, which
states that for any domain whose causal structure satisfies a sparsity condition,
a sufficiently high-dimensional learned embedding guarantees:
\textbf{(i)}~linear separability of causal features from spurious directions via
concentration of measure; \textbf{(ii)}~sub-linear sample complexity
$O(k\log(d/k))$ for recovering $k$ causal directions in a $d$-dimensional space;
\textbf{(iii)}~alignment of principal components with causal directions at
convergence, providing a theoretical foundation for PCA-based and sparse-autoencoder-based interpretability;
\textbf{(iv)}~information-theoretic optimality of superposition
(representing $\exp(\Omega(d))$ features in $d$ dimensions) under feature
sparsity; and \textbf{(v)}~a \emph{Discovery--Formalisation Separation
Theorem} proving that learning has two provably distinct phases---discovery of
causal structure in high dimensions and formalisation via compression---and that
the second phase \emph{cannot succeed without the first}: compressing before
discovering is provably suboptimal by $\Omega(\sigma^2\sqrt{k})$.  This
separation explains why post-hoc embedding compression
improves performance (the pre-trained model has already completed discovery),
why the lottery ticket hypothesis holds (the large network discovers the winning
subnetwork), and why knowledge distillation works (the teacher completes
discovery; the student only formalises).  Our framework unifies the Linear
Representation Hypothesis, the Superposition Hypothesis, scaling laws, and
benign overfitting into a single account grounded in high-dimensional
probability theory.  We validate the predictions empirically on trained
transformer embeddings from prediction markets, language models, and synthetic
data with known causal structure.
\end{abstract}

\vspace{1em}
\noindent\textbf{Keywords:}
blessing of dimensionality $\cdot$
concentration of measure $\cdot$
causal representation learning $\cdot$
superposition hypothesis $\cdot$
linear representation hypothesis $\cdot$
scaling laws $\cdot$
sparse recovery

% ══════════════════════════════════════════════════════════════════════
\section{Introduction}\label{sec:intro}
% ══════════════════════════════════════════════════════════════════════

A central puzzle of modern machine learning is that \emph{more is more}.
Classical statistical learning theory predicts that models with far more
parameters than data points should overfit catastrophically---the
\emph{curse of dimensionality}~\citep{bellman1961adaptive}.  Yet the
large-language-model revolution proceeded by violating precisely this
prescription: GPT-3 has 175 billion parameters trained on roughly 300 billion
tokens, and its successors have only grown larger, with performance consistently
improving~\citep{kaplan2020scaling, hoffmann2022training}.

This paper argues that the resolution of this puzzle lies not in any novel
mechanism but in a body of classical results from high-dimensional probability
that has been insufficiently connected to the deep learning literature: the
\textbf{blessing of dimensionality}~\citep{donoho2000high, gorban2018blessing}.

The blessing of dimensionality, a term introduced by~\citet{donoho2000high} and
developed mathematically by~\citet{gorban2018blessing}, refers to a collection of
concentration-of-measure phenomena in high-dimensional spaces.  The most striking
of these is the \emph{stochastic separation theorem}: i.i.d.\ random points
drawn from a high-dimensional distribution are, with high probability, all
\emph{linearly separable} from the rest of the set, even when the set is
exponentially large in the dimension~\citep{gorban2018blessing,
gorban2020blessing}.  In other words, high dimensions do not make classification
harder---they make it easier, provided the structure is exploited correctly.

Meanwhile, the mechanistic interpretability community has discovered that neural
networks represent far more \emph{features} (interpretable directions in
activation space) than they have neurons.  The \emph{Superposition
Hypothesis}~\citep{elhage2022superposition} proposes that networks pack
exponentially many nearly-orthogonal feature directions into their
finite-dimensional activation spaces, exploiting the geometric fact (a
consequence of the Johnson--Lindenstrauss lemma~\citep{johnson1984extensions})
that high-dimensional spaces can accommodate exponentially many
nearly-orthogonal vectors.  The \emph{Linear Representation
Hypothesis}~\citep{park2024linear} formalises the observation that high-level
concepts---gender, sentiment, factuality---are represented as linear directions
in the embedding space of large language models.

A third line of work has shown that \emph{overparameterisation} is not merely
tolerable but \emph{necessary} for good generalisation in the interpolating
regime.  The theory of \emph{benign overfitting}~\citep{bartlett2020benign}
proves that minimum-norm interpolants generalise well precisely when the data
covariance has a heavy spectral tail---that is, when there are many
``unimportant'' directions in parameter space.  The \emph{double descent}
phenomenon~\citep{belkin2019reconciling} shows that test error decreases again
once the model is sufficiently overparameterised, with the effective dimension of
the model---not the raw parameter count---governing generalisation.

These five threads---concentration of measure, superposition, linear
representations, scaling laws, and benign overfitting---have developed largely in
parallel.  This paper unifies them.

\subsection{Contributions}

We make the following contributions:

\begin{enumerate}[leftmargin=2em,itemsep=3pt]
  \item \textbf{The Universal Embedding Theorem} (Theorem~\ref{thm:main}):  
    a formal statement that any domain with $k$-sparse causal structure admits a
    learned embedding in $\R^d$ ($d \gg k$) in which causal features are
    linearly separable, recoverable from $O(k\log(d/k))$ samples, and aligned
    with the top principal components at convergence.

  \item \textbf{The PCA--Causality Alignment Theorem}
    (Theorem~\ref{thm:pca-alignment}): a novel result showing that the embedding
    gradient decomposition~\citep[Ch.~2]{omega2026} implies, at convergence, that
    PCA on the trained embedding recovers the causal directions up to an error
    controlled by the \emph{inaccessible gradient} term $\partial f/\partial
    \mathbf{X}$.

  \item \textbf{A superposition capacity bound}
    (Theorem~\ref{thm:superposition}): a quantitative version of the
    Superposition Hypothesis, showing that $d$ neurons can represent
    $\exp(\Omega(d))$ features with per-feature interference bounded by
    $O(\sqrt{\log m / d})$, and that this is information-theoretically optimal
    under sparsity.

  \item \textbf{The Discovery--Formalisation Separation Theorem}
    (Theorem~\ref{thm:discovery-formalisation}): a proof that learning has two
    provably distinct phases---\emph{discovery} (identification of causal
    structure in high dimensions) and \emph{formalisation} (compression onto the
    discovered structure)---and that formalisation without prior discovery is
    provably suboptimal.  This provides a unified explanation for post-hoc
    embedding compression, the lottery ticket hypothesis, knowledge
    distillation, and gradient-boosted tree pruning.

  \item \textbf{Empirical validation} on trained transformer embeddings from
    prediction markets (Polymarket), standard language models (GPT-2, Llama), and
    synthetic data with known causal structure, including a counterintuitive
    prediction that adding random noise dimensions to an embedding
    \emph{improves} causal feature recovery.
\end{enumerate}


\subsection{Outline}

Section~\ref{sec:framework} establishes the data pipeline and embedding
framework.  Section~\ref{sec:blessing} reviews the mathematical foundations of
the blessing of dimensionality.  Section~\ref{sec:main-results} states and proves
the main theorems, including the Discovery--Formalisation Separation.
Section~\ref{sec:connections} draws the connections to scaling laws, double
descent, and the interpretability literature.
Section~\ref{sec:experiments} presents empirical validation.
Section~\ref{sec:discussion} discusses implications and limitations.


% ══════════════════════════════════════════════════════════════════════
\section{The Learning Framework}\label{sec:framework}
% ══════════════════════════════════════════════════════════════════════

We adopt the data-pipeline formalism of~\citet{omega2026}, which separates
representation learning from prediction in a way that maps naturally onto
transformer architectures.

\begin{definition}[Data pipeline]\label{def:pipeline}
A \emph{learning pipeline} consists of:
\begin{enumerate}[label=(\roman*),leftmargin=2.5em,itemsep=2pt]
  \item An \emph{observation space} $\cO$ and an observation $\cX \in \cO$;
  \item A \emph{proposition map} $\tau : \cO \to \Pi$, producing a structured
    representation $X = \tau(\cX)$ in a proposition space~$\Pi$;
  \item A \emph{parametric embedding} $\phi_\psi : \Pi \to \R^d$ with parameters
    $\psi \in \R^p$, producing a feature vector
    $\mathbf{X} = \phi_\psi(X) \in \R^d$;
  \item A \emph{true causal map} $f : \R^d \to \R^q$ (unknown), defining the
    target $Y = f(\mathbf{X})$;
  \item A \emph{parametric predictor} $\hat{f}_\theta : \R^d \to \R^q$ with
    parameters $\theta \in \R^l$, producing predictions
    $\hat{Y} = \hat{f}_\theta(\mathbf{X})$;
  \item A \emph{loss function}
    $\cL : \R^q \times \R^q \to \R$, continuously differentiable, with
    $\cL(Y,\hat{Y}) = \cL\bigl(f(\mathbf{X}_\psi),\,
    \hat{f}_\theta(\mathbf{X}_\psi)\bigr)$.
\end{enumerate}
We write $\beta = (\theta,\psi) \in \R^{l+p}$ for the full parameter vector.
\end{definition}

This decomposition is natural in modern architectures: in a transformer, $\psi$
are the embedding and encoder parameters and $\theta$ are the prediction-head
parameters.  In a gradient-boosted model with engineered features, $\psi$ encodes
the feature construction and $\theta$ the tree ensemble.

\subsection{The Gradient Decomposition}\label{sec:grad-decomp}

The embedding--predictor separation produces structurally different gradients for
each parameter group.

\begin{proposition}[Predictor gradient]\label{prop:pred-grad}
Since $Y = f(\mathbf{X}_\psi)$ does not depend on~$\theta$:
\begin{equation}\label{eq:pred-grad}
  \nabla_\theta \cL
  = \frac{\partial \cL}{\partial \hat{Y}}
    \cdot \frac{\partial \hat{f}_\theta(\mathbf{X}_\psi)}{\partial \theta}
  \;\in\; \R^{1\times l}.
\end{equation}
\end{proposition}

\begin{proposition}[Embedding gradient]\label{prop:embed-grad}
Both arguments of $\cL$ depend on $\psi$ through
$\mathbf{X}_\psi = \phi_\psi(X)$.  By the multivariate chain rule:
\begin{equation}\label{eq:embed-grad}
  \nabla_\psi \cL
  = \underbrace{
      \frac{\partial \cL}{\partial Y}
      \cdot \frac{\partial f}{\partial \mathbf{X}}
      \cdot \frac{\partial \phi_\psi}{\partial \psi}
    }_{\text{through } f \text{ (inaccessible)}}
  + \underbrace{
      \frac{\partial \cL}{\partial \hat{Y}}
      \cdot \frac{\partial \hat{f}_\theta}{\partial \mathbf{X}}
      \cdot \frac{\partial \phi_\psi}{\partial \psi}
    }_{\text{through } \hat{f}_\theta \text{ (computable)}}
  \;\in\; \R^{1\times p}.
\end{equation}
\end{proposition}

\begin{remark}[The inaccessible term]\label{rmk:inaccessible}
The first term in~\eqref{eq:embed-grad} requires
$\partial f/\partial \mathbf{X}$---the Jacobian of the true causal map.  When $f$
is unknown, this term is \emph{inaccessible}: backpropagation computes only the
second term.  The resulting bias vanishes only when
$\hat{f}_\theta = f$ (perfect prediction) or when
$\partial \cL / \partial Y = 0$.  This inaccessible term will play a crucial role
in bounding the PCA--causality alignment error
(Theorem~\ref{thm:pca-alignment}).
\end{remark}


\subsection{Causal Sparsity}\label{sec:sparsity}

The structural assumption that makes the blessing of dimensionality operative:

\begin{assumption}[Causal $k$-sparsity]\label{ass:sparsity}
The true causal map $f:\R^d \to \R^q$ depends on $\mathbf{X}$ only through its
projection onto a $k$-dimensional subspace $V_f \subset \R^d$, with $k \ll d$.
Formally, there exists an orthogonal projector $P_f \in \R^{d\times d}$ with
$\rank(P_f) = k$ such that $f(\mathbf{X}) = f(P_f \mathbf{X})$ for all
$\mathbf{X}$.
\end{assumption}

The assumption is well-motivated across domains:

\begin{example}[Language]
Of the $\sim$50{,}000 token directions in a GPT-2 embedding space ($d = 768$),
any given next-token prediction depends on a handful of semantic features
(topic, sentiment, syntactic role, a few entity-specific
directions)~\citep{park2024linear, bricken2023monosemanticity}.  Empirically,
sparse autoencoders find that most tokens activate only $\sim$10--50 features
out of millions available~\citep{templeton2024scaling}.
\end{example}

\begin{example}[Financial markets]
Of the hundreds of potential drivers of a prediction-market contract's price,
typically 5--15 are causally relevant at any time: a few political events,
macroeconomic indicators, and information-flow sources.
\end{example}

\begin{example}[Biological systems]
Of the $\sim$20{,}000 genes in the human genome, any given phenotype is driven
by a sparse set of causal pathways~\citep{boyle2017expanded}.
\end{example}


% ══════════════════════════════════════════════════════════════════════
\section{The Blessing of Dimensionality: Mathematical Foundations}%
\label{sec:blessing}
% ══════════════════════════════════════════════════════════════════════

We review the key results from high-dimensional probability that underpin the
main theorems.  The presentation follows~\citet{gorban2018blessing}
and~\citet{vershynin2018high}.

\subsection{Concentration of Measure}\label{sec:concentration}

The foundational phenomenon is that Lipschitz functions of many independent (or
weakly dependent) random variables concentrate sharply around their mean.

\begin{theorem}[Lévy's lemma; see {\citealt[Ch.~5]{vershynin2018high}}]%
\label{thm:levy}
Let $f : \mathbb{S}^{d-1} \to \R$ be $L$-Lipschitz with respect to geodesic
distance, and let $\mathbf{x}$ be drawn uniformly from the unit sphere
$\mathbb{S}^{d-1}$.  Then for all $t > 0$:
\begin{equation}\label{eq:levy}
  \Prob\bigl[\abs{f(\mathbf{x}) - \E[f(\mathbf{x})]} \geq t\bigr]
  \;\leq\; 2\exp\!\Bigl(-\frac{(d-1)\,t^2}{2L^2}\Bigr).
\end{equation}
\end{theorem}

This implies that in high dimensions, the squared distance from any fixed point
to a random point concentrates sharply around its mean---the \emph{thin shell}
phenomenon.  All random points land on a thin spherical shell of predictable
radius.

\subsection{Stochastic Separation Theorems}\label{sec:separation}

The stochastic separation theorems of~\citet{gorban2018blessing} describe the
\emph{fine structure} of these thin shells:

\begin{theorem}[Stochastic separation; {\citealt[Thm.~1]{gorban2018blessing}}]%
\label{thm:stochastic-sep}
Let $\mathbf{x}_1, \ldots, \mathbf{x}_M$ be i.i.d.\ random points from a
distribution in $\R^d$ with effective dimension $\deff$.  For each $i$, the
probability that $\mathbf{x}_i$ can be linearly separated from
$\{\mathbf{x}_j\}_{j \neq i}$ by Fisher's linear discriminant satisfies:
\begin{equation}\label{eq:sep}
  \Prob[\text{$\mathbf{x}_i$ is linearly separable}]
  \;\geq\; 1 - M \cdot \exp(-C\,\deff),
\end{equation}
where $C > 0$ is a universal constant.  In particular, even exponentially large
sets ($M = \exp(\alpha\, \deff)$ for $\alpha < C$) are fully linearly separable
with high probability.
\end{theorem}

The effective dimension $\deff$ (not the ambient dimension $d$) governs
separability.  We make this precise:

\begin{definition}[Effective dimension]\label{def:deff}
For a random vector $\mathbf{X} \in \R^d$ with covariance $\Sigma$, the
\emph{effective dimension} is:
\begin{equation}\label{eq:deff}
  \deff(\Sigma) \;=\; \frac{\tr(\Sigma)}{\lambda_{\max}(\Sigma)}
  \;=\; \frac{\sum_{i=1}^{d} \lambda_i}{\lambda_1},
\end{equation}
where $\lambda_1 \geq \cdots \geq \lambda_d \geq 0$ are the eigenvalues of
$\Sigma$.  If only $k$ eigenvalues are significant, $\deff \approx k$.
\end{definition}

\begin{remark}[Triple-duty effective dimension]\label{rmk:triple-duty}
The quantity $\deff$ simultaneously governs three distinct aspects of the
embedding, a unification that is not a design choice but a mathematical
consequence~\citep[cf.\ the triple-duty evidence weight, \S10.3]{omega2026}:
\begin{enumerate}[leftmargin=2em,itemsep=1pt]
  \item \textbf{Gradient concentration quality:} The sample complexity for
    recovering the true gradient direction is
    $M = O(\deff\log d / \varepsilon^2)$, not $O(d\log d/\varepsilon^2)$
    (Theorem~\ref{thm:main}(ii)).
  \item \textbf{Intrinsic causal dimension:} When the causal map is $k$-sparse,
    $\deff \approx k$ reveals the true complexity of the prediction task,
    independent of the ambient dimension $d$
    (Theorem~\ref{thm:pca-alignment}).
  \item \textbf{Compressibility:} The ratio $\deff/d$ measures how much the
    embedding can be compressed without loss of causal information---the
    achievable compression factor is $d/\deff$
    (Theorem~\ref{thm:discovery-formalisation}(b)).
\end{enumerate}
All three are manifestations of the spectral structure of the embedding
covariance $\Sigma$: the eigenvalue gap between the $k$-th and $(k{+}1)$-th
eigenvalues simultaneously creates gradient concentration, identifies the causal
subspace, and defines the compression boundary.
\end{remark}

\subsection{The Johnson--Lindenstrauss Lemma}\label{sec:jl}

\begin{theorem}[Johnson--Lindenstrauss; {\citealt{johnson1984extensions}}]%
\label{thm:jl}
For any $0 < \varepsilon < 1$, any integer $n \geq 1$, and any set
$S \subset \R^D$ with $|S| = n$, there exists a linear map
$A : \R^D \to \R^d$ with
$d = O(\varepsilon^{-2}\log n)$
such that for all $\mathbf{u}, \mathbf{v} \in S$:
\begin{equation}\label{eq:jl}
  (1-\varepsilon)\norm{\mathbf{u}-\mathbf{v}}^2
  \;\leq\; \norm{A\mathbf{u} - A\mathbf{v}}^2
  \;\leq\; (1+\varepsilon)\norm{\mathbf{u}-\mathbf{v}}^2.
\end{equation}
\end{theorem}

\begin{remark}[JL and superposition]
The converse reading is equally important: $d$ dimensions can accommodate
$n = \exp(\Omega(\varepsilon^2 d))$ points (or directions) whose pairwise
distances (or angles) are preserved to within $(1\pm\varepsilon)$.  This is the
geometric foundation of the Superposition Hypothesis~\citep{elhage2022superposition}:
a $d$-dimensional activation space can represent
$m = \exp(\Omega(d))$ nearly-orthogonal feature directions.
\end{remark}


\subsection{Sparse Recovery and the Restricted Isometry Property}%
\label{sec:rip}

\begin{theorem}[Candès--Tao sparse recovery; {\citealt{candes2005decoding}}]%
\label{thm:rip}
Let $\mathbf{g} \in \R^d$ be $k$-sparse (i.e.\ $\norm{\mathbf{g}}_0 \leq k$)
and let $\Phi \in \R^{M\times d}$ be a measurement matrix satisfying the
\emph{Restricted Isometry Property} of order $2k$ with constant
$\delta_{2k} < \sqrt{2}-1$.  Then $\mathbf{g}$ can be exactly recovered from
$\mathbf{y} = \Phi\,\mathbf{g}$ via $\ell_1$-minimisation:
\begin{equation}\label{eq:l1-recovery}
  \hat{\mathbf{g}} = \argmin_{\mathbf{z}} \norm{\mathbf{z}}_1
  \quad\text{subject to}\quad \Phi\,\mathbf{z} = \mathbf{y},
\end{equation}
provided $M \geq C \cdot k\log(d/k)$.
\end{theorem}

The critical observation: going from $d = 10$ to $d = 10{,}000$ action
dimensions costs an additional factor of $\log(1000) \approx 7$ samples, not
$1000$.  This is the mechanism by which the curse is transmuted into a blessing
under sparsity.


% ══════════════════════════════════════════════════════════════════════
\section{Main Results}\label{sec:main-results}
% ══════════════════════════════════════════════════════════════════════

We now state and prove the main theorems.  Throughout, we work with a data
pipeline (Definition~\ref{def:pipeline}) satisfying the causal $k$-sparsity
assumption (Assumption~\ref{ass:sparsity}).

\subsection{The Universal Embedding Theorem}\label{sec:main-thm}

\begin{theorem}[Universal Embedding Theorem]\label{thm:main}
Let the data pipeline (Definition~\ref{def:pipeline}) satisfy
Assumption~\ref{ass:sparsity} with causal subspace $V_f$ of dimension $k$.  Let
$\phi_\psi : \Pi \to \R^d$ be a learned embedding with $d \geq C_0\, k\log(k/\delta)$
for a universal constant $C_0$.  Suppose the embedding distribution
$\mathbf{X} = \phi_\psi(X)$ is sub-Gaussian with effective dimension
$\deff(\Sigma_{\mathbf{X}}) \geq c_0\,d$ for some $c_0 > 0$.  Then with
probability at least $1-\delta$:

\begin{enumerate}[label=\textup{(\roman*)},leftmargin=2.5em,itemsep=4pt]
  \item \textbf{Linear separability.}
    Any causal feature direction $\mathbf{v} \in V_f$ can be linearly separated
    from any set of $M \leq \exp(\alpha\,\deff)$ spurious directions by
    Fisher's linear discriminant, for a constant $\alpha$ depending only on the
    sub-Gaussian norm.

  \item \textbf{Sub-linear recovery.}
    The projection of the true gradient $\nabla_{\mathbf{X}}\cL$ onto the causal
    subspace $V_f$ can be recovered from
    \begin{equation}\label{eq:sublinear}
      M \;=\; O\!\left(\frac{k\log(d/k)}{\varepsilon^2}\right)
    \end{equation}
    i.i.d.\ stochastic gradient samples, to $\varepsilon$-accuracy in $\ell_2$
    norm.

  \item \textbf{Superposition capacity.}
    The embedding space $\R^d$ can simultaneously represent
    $m = \exp(\Omega(\varepsilon^2 d))$ feature directions such that all pairwise
    inner products satisfy $|\ip{\mathbf{v}_i}{\mathbf{v}_j}| \leq \varepsilon$
    for $i\neq j$.

  \item \textbf{PCA--causality alignment} (at convergence).
    If the predictor has converged to $\hat{f}_\theta = f + \eta$ with
    $\norm{\eta}_{\mathrm{op}} \leq \varepsilon_f$, then the top-$k$ principal
    components of the embedding covariance $\Sigma_\mathbf{X}$ span a subspace
    $\hat{V}_k$ satisfying:
    \begin{equation}\label{eq:alignment}
      \norm{P_{V_f} - P_{\hat{V}_k}}_F
      \;\leq\; C_1\,\varepsilon_f \cdot
      \frac{\lambda_{k+1}(\Sigma)}{\lambda_k(\Sigma) - \lambda_{k+1}(\Sigma)},
    \end{equation}
    where $P_{V_f}$, $P_{\hat{V}_k}$ are the orthogonal projectors onto $V_f$
    and $\hat{V}_k$ respectively.
\end{enumerate}
\end{theorem}

Parts (i)--(iii) follow from the classical results of
Sections~\ref{sec:separation}--\ref{sec:rip} applied to the embedding space.
Part~(iv) is the principal novel contribution and requires the embedding gradient
decomposition.

\subsection{Proof of Parts (i)--(iii)}\label{sec:proof-i-iii}

\begin{proof}[Proof of (i)]
Apply Theorem~\ref{thm:stochastic-sep} to the embedding distribution
$\mathbf{X} = \phi_\psi(X) \in \R^d$.  The sub-Gaussianity assumption ensures
that the concentration inequalities hold with the stated constants.  The causal
direction $\mathbf{v} \in V_f$ is one point in $\R^d$; the spurious directions
are $M$ additional points.  By~\eqref{eq:sep}, linear separation holds with
probability $\geq 1 - M\exp(-C\,\deff)$.  For $M = \exp(\alpha\,\deff)$ with
$\alpha < C$, this is $\geq 1 - \exp(-(\alpha - C)\,\deff)$.
\end{proof}

\begin{proof}[Proof of (ii)]
The causal $k$-sparsity assumption (Assumption~\ref{ass:sparsity}) implies that
the true gradient $\nabla_{\mathbf{X}} f$ has column space contained in $V_f$,
hence rank at most~$k$.  The gradient of the loss with respect to the embedding,
$\nabla_{\mathbf{X}} \cL$, inherits this sparsity structure:
\[
  \nabla_{\mathbf{X}} \cL
  = \frac{\partial\cL}{\partial Y}\cdot \frac{\partial f}{\partial \mathbf{X}}
  \in V_f.
\]
Each stochastic gradient sample $\hat{\mathbf{g}}_m$ provides a noisy measurement
of this $k$-sparse signal.  By Theorem~\ref{thm:rip}, the sample-averaged
gradient concentrates around the true gradient with
$M = O(k\log(d/k)/\varepsilon^2)$ samples.  More precisely, define the
measurement matrix $\Phi \in \R^{M\times d}$ whose rows are the stochastic
gradient samples (centred and normalised).  Standard results on sub-Gaussian
random matrices~\citep[Thm.~9.1.1]{vershynin2018high} show that $\Phi$ satisfies
the RIP of order $2k$ with high probability for the stated $M$.
\end{proof}

\begin{proof}[Proof of (iii)]
This is a direct application of Theorem~\ref{thm:jl}.  We require $m$ unit
vectors $\mathbf{v}_1,\ldots,\mathbf{v}_m \in \R^d$ with
$|\ip{\mathbf{v}_i}{\mathbf{v}_j}| \leq \varepsilon$.  By the JL lemma, a
random $d$-dimensional subspace of $\R^m$ preserves all pairwise inner products
to within $\varepsilon$ provided $d = \Omega(\varepsilon^{-2}\log m)$, hence
$m = \exp(\Omega(\varepsilon^2 d))$.  This is the quantitative content of the
Superposition Hypothesis~\citep{elhage2022superposition}.
\end{proof}


\subsection{The PCA--Causality Alignment Theorem}\label{sec:pca-thm}

This is the principal novel contribution.

\begin{theorem}[PCA--Causality Alignment]\label{thm:pca-alignment}
Adopt the setup of Theorem~\ref{thm:main}.  Let the predictor have converged to
$\hat{f}_\theta$ with uniform approximation error
$\norm{\hat{f}_\theta - f}_\infty \leq \varepsilon_f$.  Let $\Sigma_\psi^*$
denote the covariance of the embedding $\phi_{\psi^*}(X)$ at the converged
parameters $\psi^*$, and let $\hat{V}_k$ be the span of its top-$k$
eigenvectors.  Then:
\begin{equation}\label{eq:pca-main}
  \sin\angle(V_f,\,\hat{V}_k)
  \;\leq\;
  \frac{C_2\,\varepsilon_f\cdot\norm{\nabla_{\mathbf{X}} f}_{\mathrm{op}}}
       {\lambda_k(\Sigma_\psi^*) - \lambda_{k+1}(\Sigma_\psi^*)},
\end{equation}
where $\angle(V_f, \hat{V}_k)$ is the largest principal angle between the two
$k$-dimensional subspaces.
\end{theorem}

\begin{proof}[Proof sketch]
The proof proceeds in three steps.

\textbf{Step 1: Gradient covariance structure.}
At convergence, the gradient signal that shaped the embedding during training was
dominated by the computable term in~\eqref{eq:embed-grad}.  Let
$G_\psi = \E[\nabla_\psi \cL \cdot (\nabla_\psi \cL)^\top]$ be the gradient
covariance.  The computable term has factor
$(\partial\hat{f}_\theta/\partial\mathbf{X})\cdot(\partial\phi_\psi/\partial\psi)$,
whose column space converges to $V_f$ as $\hat{f}_\theta \to f$.  The
inaccessible term contributes only to directions in $V_f$ (by the sparsity
assumption), so it reinforces rather than perturbs the alignment.

\textbf{Step 2: Implicit regularisation and embedding covariance.}
Gradient descent with learning rate $\gamma$ on the embedding parameters $\psi$
induces an implicit covariance structure on the embedding: the directions that
received the largest cumulative gradient signal acquire the largest variance.
Under standard assumptions on the learning dynamics (see
Appendix~\ref{app:implicit-reg}), the converged embedding covariance satisfies:
\begin{equation}\label{eq:sigma-alignment}
  \Sigma_\psi^* \;\approx\;
  P_{V_f}\,\Sigma_\psi^*\,P_{V_f} \;+\; \sigma_\perp^2\,(I - P_{V_f}),
\end{equation}
where $\sigma_\perp^2$ is the residual variance in the orthogonal complement of
$V_f$, driven by the inaccessible gradient bias and the noise in the training
data.

\textbf{Step 3: Davis--Kahan perturbation.}
The decomposition~\eqref{eq:sigma-alignment} shows that $\Sigma_\psi^*$ is a
rank-$k$ perturbation of $\sigma_\perp^2 I$.  By the Davis--Kahan
$\sin\theta$ theorem~\citep{davis1970rotation}, the angle between the top-$k$
eigenspace of $\Sigma_\psi^*$ and $V_f$ is bounded by:
\[
  \sin\angle(V_f,\,\hat{V}_k)
  \;\leq\;
  \frac{\norm{\Sigma_\psi^* - (\tilde{\Sigma}_f + \sigma_\perp^2 I)}_{\mathrm{op}}}
       {\lambda_k(\Sigma_\psi^*) - \lambda_{k+1}(\Sigma_\psi^*)},
\]
where $\tilde{\Sigma}_f = P_{V_f}\,\Sigma_\psi^*\,P_{V_f}$ is the ``ideal''
covariance.  The numerator is controlled by $\varepsilon_f$ (the predictor
approximation error), giving the claimed bound.
\end{proof}

\begin{remark}[The spectral gap condition]
The denominator $\lambda_k - \lambda_{k+1}$ requires a spectral gap between the
causal and non-causal eigenvalues.  This is precisely the condition that the
$k$~causal features are ``more important'' than all non-causal features---a
condition that the blessing of dimensionality helps satisfy, because
concentration of measure ensures that noise variance distributes uniformly across
the $d - k$ non-causal directions, keeping $\lambda_{k+1}$ small when $d$ is
large.
\end{remark}


\subsection{Superposition Capacity and Optimality}\label{sec:superposition-thm}

\begin{theorem}[Superposition is optimal under sparsity]%
\label{thm:superposition}
Let a representation system have $d$ neurons and represent $m$ Boolean features,
of which at most $s$ are simultaneously active (feature sparsity).  Then:
\begin{enumerate}[label=\textup{(\alph*)},leftmargin=2.5em]
  \item \textbf{Capacity:} There exist $m = \exp(\Omega(d))$ unit vectors
    $\mathbf{v}_1,\ldots,\mathbf{v}_m \in \R^d$ such that for any subset $S$
    with $|S| \leq s$, the active features $\{\mathbf{v}_i\}_{i\in S}$ can be
    decoded with per-feature error at most
    $O(s\sqrt{\log m / d})$.

  \item \textbf{Optimality:} Any scheme representing $m$ features in $d$
    dimensions with $s$-sparse activation requires per-feature error at least
    $\Omega(s/\sqrt{d})$.  The superposition scheme achieves this up to a
    $\sqrt{\log m}$ factor.

  \item \textbf{Computation gap:}
    While passive \emph{representation} can encode $\exp(\Omega(d))$ features,
    active \emph{computation} on $m'$ output features in superposition
    requires $\Omega(m'\log^2 m')$ parameters~\citep{adler2024complexity},
    implying an exponential gap between representational and computational
    capacity.
\end{enumerate}
\end{theorem}

\begin{proof}[Proof of (a)]
Construct the $m$ feature vectors by drawing i.i.d.\ rows of a random matrix
$V \in \R^{m\times d}$ with entries $V_{ij} \sim \cN(0,1/d)$.  For any pair
$i\neq j$:
\[
  \E[\ip{\mathbf{v}_i}{\mathbf{v}_j}] = 0, \qquad
  \mathrm{Var}(\ip{\mathbf{v}_i}{\mathbf{v}_j}) = 1/d.
\]
By the Hanson--Wright inequality, $|\ip{\mathbf{v}_i}{\mathbf{v}_j}| \leq
t/\sqrt{d}$ with probability $\geq 1 - 2\exp(-ct^2)$.  Taking a union bound
over all $\binom{m}{2}$ pairs and setting $t = C\sqrt{\log m}$, all pairwise
inner products are bounded by $C\sqrt{\log m / d}$ with probability
$\geq 1 - 2m^{-c'}$, provided $d = \Omega(\log m)$.

Given an activation $\mathbf{a} = \sum_{i\in S} \alpha_i \mathbf{v}_i$ with
$|S| = s$, decoding feature $j$ via $\hat{\alpha}_j = \ip{\mathbf{v}_j}{\mathbf{a}}$
incurs interference:
\[
  \hat{\alpha}_j - \alpha_j = \sum_{i\in S\setminus\{j\}} \alpha_i\,
  \ip{\mathbf{v}_j}{\mathbf{v}_i}.
\]
The error is bounded by
$s \cdot \max_{i\neq j}|\ip{\mathbf{v}_j}{\mathbf{v}_i}| =
O(s\sqrt{\log m / d})$.
\end{proof}


\subsection{Effective Evidence Weight and the Natural Curriculum}%
\label{sec:evidence}

Combining the sparsity-aware dimension with evidence weighting yields an
adaptive training dynamic.

\begin{definition}[Effective evidence weight]
For agent (or data source) $i$ with gradient covariance $\Sigma_i$, define:
\begin{equation}\label{eq:eff-evidence}
  V_i^{\mathrm{eff}} = \frac{\deff(\Sigma_i)}{d}\cdot V_i,
  \qquad
  w_i^{\mathrm{eff}} = \frac{V_{\min}^{\mathrm{eff}}}{V_i^{\mathrm{eff}}},
\end{equation}
where $V_i = \tr(\Sigma_i)$ is the total gradient variance.
\end{definition}

\begin{proposition}[Evidence amplification]\label{prop:evidence-amp}
For an agent with $k$-sparse optimal policy in a $d$-action space:
\begin{equation}\label{eq:evidence-amp}
  \frac{w_i^{\mathrm{eff}}}{w_i} = \frac{d}{\deff(\Sigma_i)} \geq \frac{d}{k}.
\end{equation}
For $d = 10{,}000$ and $k = 5$: the effective weight is $2{,}000\times$ larger.
\end{proposition}

This amplification creates a natural curriculum:
\begin{enumerate}[leftmargin=2em]
  \item \textbf{Early training:} $\ell_1$ regularisation drives extreme sparsity
    ($k \approx 1$--$2$), $\deff \approx k$, high effective evidence, fast
    convergence on the dominant causal features.
  \item \textbf{Mid training:} Support grows, but evidence improves
    simultaneously.  The model explores new features with calibrated confidence.
  \item \textbf{Late training:} $\ell_1$ anneals to zero; full support recovered.
    Convergence to the true (possibly dense) optimum.
\end{enumerate}


\subsection{The Discovery--Formalisation Separation}%
\label{sec:discovery-formalisation}

The preceding results establish that high-dimensional embeddings \emph{can}
recover causal structure.  We now prove a much stronger claim: high-dimensional
discovery is \emph{necessary}---it cannot be bypassed by learning directly in
low dimensions.  Learning has two provably distinct phases, and the second
cannot succeed without the first.

\subsubsection{The Two Phases of Learning}

The distinction mirrors a classical idea from the philosophy of mathematics.
P\'olya's \emph{Patterns of Plausible Inference}~\citep{polya1954patterns}
distinguishes \emph{heuristic reasoning}---exploring a large space of
possibilities to identify plausible conjectures---from \emph{demonstrative
reasoning}---compressing the discovered conjecture into a parsimonious proof.
The same separation arises in gradient-boosted decision trees: the tree is
\emph{overgrown} to discover all relevant split points, then \emph{pruned} to
retain only those that generalise.  We now prove that this pattern is not merely
a heuristic but a mathematical necessity.

\begin{definition}[Discovery and formalisation]\label{def:disc-form}
Let $\phi_\psi^{(d)} : \Pi \to \R^d$ be a $d$-dimensional embedding and
$\pi_k : \R^d \to \R^k$ be a $k$-dimensional compression (e.g., PCA projection
onto the top-$k$ components).  We define:
\begin{itemize}[leftmargin=2em]
  \item \textbf{Discovery:} Training $\phi_\psi^{(d)}$ in $\R^d$ with
    $d \geq C_0\,k\log(k/\delta)$.  The output is the converged embedding and
    its spectral structure.
  \item \textbf{Formalisation:} Applying $\pi_k$ to the converged embedding to
    obtain a $k$-dimensional representation
    $\tilde{\mathbf{X}} = \pi_k(\phi_{\psi^*}^{(d)}(X)) \in \R^k$.
  \item \textbf{Direct learning:} Training $\phi_\psi^{(k)} : \Pi \to \R^k$
    directly in $k$ dimensions, bypassing the high-dimensional phase.
\end{itemize}
\end{definition}

\begin{theorem}[Discovery--Formalisation Separation]%
\label{thm:discovery-formalisation}
Let the data pipeline satisfy Assumption~\ref{ass:sparsity} with causal
dimension~$k$, and let $\varepsilon$ be label noise with variance~$\sigma^2$.
Then:

\begin{enumerate}[label=\textup{(\alph*)},leftmargin=2.5em,itemsep=4pt]
  \item \textbf{Discovery requires high dimensions.}
    For any embedding $\phi_\psi : \Pi \to \R^d$ with $d < k$, there exist
    data distributions satisfying Assumption~\ref{ass:sparsity} such that no
    estimator based on $\phi_\psi(X)$ can recover the causal subspace $V_f$ with
    error less than a constant:
    \begin{equation}\label{eq:disc-lower}
      \inf_{\hat{V}_k} \;
      \E\bigl[\sin\angle(V_f, \hat{V}_k)\bigr]
      \;\geq\; c_0 > 0,
    \end{equation}
    where the infimum is over all estimators $\hat{V}_k$ of $V_f$ based on
    $n$ observations of $(\phi_\psi(X), Y)$.

  \item \textbf{Formalisation after discovery is optimal.}
    Let $\phi_{\psi^*}^{(d)}$ be the converged embedding in $\R^d$
    ($d \geq C_0 k\log(k/\delta)$) and
    $\pi_k : \R^d \to \R^k$ be the PCA projection onto its top-$k$ eigenspace.
    Then the compressed predictor achieves risk:
    \begin{equation}\label{eq:formal-risk}
      \mathrm{Risk}(\hat{f}_\theta \circ \pi_k)
      \;\leq\;
      \mathrm{Risk}(\hat{f}_\theta)
      + \sigma^2\frac{k}{n}
      + O(\varepsilon_f^2),
    \end{equation}
    where the second term is the estimation error in the $k$-dimensional
    compressed space and $\varepsilon_f$ is the approximation error from
    Theorem~\ref{thm:pca-alignment}.  When $\sigma^2 > 0$ and $d \gg k$,
    this is \emph{strictly better} than the uncompressed risk, because the
    uncompressed predictor incurs estimation error $\sigma^2 d/n \gg
    \sigma^2 k/n$.

  \item \textbf{Formalisation without discovery is suboptimal.}
    Let $\phi_{\psi^*}^{(k)} : \Pi \to \R^k$ be an embedding trained directly
    in $k$ dimensions.  Then there exist data distributions satisfying
    Assumption~\ref{ass:sparsity} such that:
    \begin{equation}\label{eq:skip-penalty}
      \mathrm{Risk}(\hat{f}_\theta \circ \phi_{\psi^*}^{(k)})
      \;\geq\;
      \mathrm{Risk}(\hat{f}_\theta \circ \pi_k \circ \phi_{\psi^*}^{(d)})
      \;+\; \Omega\!\left(\sigma^2\sqrt{k}\right).
    \end{equation}
    The direct low-dimensional embedding incurs an excess risk penalty of order
    $\sigma^2\sqrt{k}$ relative to the discover-then-formalise strategy.
\end{enumerate}
\end{theorem}

\begin{proof}[Proof of (a)]
If $d < k$, the embedding $\phi_\psi(X) \in \R^d$ maps $\R^D$ (the input
space) into a space of dimension less than $k$.  The causal subspace $V_f$
is $k$-dimensional, so the embedding must collapse at least one causal
direction.  Formally, $\phi_\psi$ composed with any linear map $\R^d \to \R^k$
has rank at most $d < k$, so it cannot span $V_f$.  By the rank--nullity
theorem, the kernel of the projection $V_f \to \R^d$ is at least
$(k-d)$-dimensional.  The collapsed causal directions contribute an
irreducible error bounded below by the variance of $Y$ along those directions,
which is $\geq c_0 > 0$ for generic $f$.
\end{proof}

\begin{proof}[Proof of (b)]
By Theorem~\ref{thm:pca-alignment}, the PCA projection $\pi_k$ onto the top-$k$
eigenspace of the converged embedding satisfies
$\sin\angle(V_f, \hat{V}_k) \leq C_2\,\varepsilon_f /
(\lambda_k - \lambda_{k+1})$.  The compressed representation
$\tilde{\mathbf{X}} = \pi_k(\mathbf{X}) \in \R^k$ retains the causal
information (up to the alignment error) while discarding the $d-k$ non-causal
dimensions.

We apply the binary risk partition~\citep[Prop.~6]{omega2026} with $W = V_f$
(causal subspace) and $S = V_f^\perp$ (non-causal complement).  The risk of any
predictor decomposes into six terms---three error sources (irreducible, bias,
variance) crossed with two subspaces ($W$, $S$):
\begin{equation}\label{eq:six-way}
  \mathrm{Risk} =
  \underbrace{R^{\mathrm{irr}}_W + R^{\mathrm{irr}}_S}_{\text{irreducible}}
  + \underbrace{R^{\mathrm{bias}}_W + R^{\mathrm{bias}}_S}_{\text{bias}}
  + \underbrace{R^{\mathrm{var}}_W + R^{\mathrm{var}}_S}_{\text{variance}}.
\end{equation}
By Assumption~\ref{ass:sparsity}, $R^{\mathrm{irr}}_S = 0$ (non-causal
directions do not affect~$Y$) and $R^{\mathrm{bias}}_S = 0$ (no target to
approximate in~$S$).  The compression step $\pi_k$ zeroes out
$R^{\mathrm{var}}_S$ (the discarded dimensions contribute no estimation noise
to the compressed predictor), leaving:
\begin{align}
  \mathrm{Risk}(\hat{f}_\theta \circ \pi_k)
  &= R^{\mathrm{irr}}_W
   + \underbrace{R^{\mathrm{bias}}_W}_{\leq\, O(\varepsilon_f^2)}
   + \underbrace{R^{\mathrm{var}}_W}_{=\, \sigma^2 k/n}
  \notag\\
  &\leq R^{\mathrm{irr}}_W + O(\varepsilon_f^2) + \sigma^2\frac{k}{n}.
  \label{eq:compressed-risk}
\end{align}
The bias $R^{\mathrm{bias}}_W = O(\varepsilon_f^2)$ because the PCA projection
misaligns with $V_f$ by at most $\sin\angle = O(\varepsilon_f)$
(Theorem~\ref{thm:pca-alignment}).  The variance $R^{\mathrm{var}}_W =
\sigma^2 k/n$ because estimation in the $k$-dimensional projected space has
effective degrees of freedom~$k$.

Compare to the \emph{uncompressed} risk, where $R^{\mathrm{var}}_S =
\sigma^2(d-k)/n$ is retained:
\begin{equation}\label{eq:uncompressed-risk}
  \mathrm{Risk}(\hat{f}_\theta)
  \geq R^{\mathrm{irr}}_W + \sigma^2\frac{k}{n}
  + \sigma^2\frac{d-k}{n}.
\end{equation}
Compression eliminates the third term entirely, yielding a variance improvement
of $\sigma^2(d-k)/n$ at a bias cost of $O(\varepsilon_f^2)$.  When $d \gg k$
and $\varepsilon_f$ is small (good predictor), the net effect is strictly
positive.  This is precisely the mechanism observed
by~\citet{drinkall2025dimensionality}: post-hoc compression of LLM embeddings
improves performance on tasks where the causal signal is sparse ($k \ll d$) and
the noise $\sigma^2$ is significant.
\end{proof}

\begin{proof}[Proof of (c)]
The key mechanism is that in $k$ dimensions, the embedding must
\emph{simultaneously} represent the $k$ causal directions \emph{and} separate
them from noise---but there is no room left for the noise to ``go.''

Formally, consider the decomposition $Y = f(P_{V_f}\mathbf{X}) + \varepsilon$.
In $\R^d$ with $d \gg k$, the noise $\varepsilon$ is approximately orthogonal to
$V_f$: the projection of $\varepsilon$ onto $V_f$ has magnitude
$O(\sigma\sqrt{k/d})$, which vanishes as $d \to \infty$.  This is the
concentration-of-measure blessing---noise spreads uniformly across all $d$
directions, so its projection onto any fixed $k$-dimensional subspace is small.

In $\R^k$, this blessing is absent.  The embedding $\phi_\psi^{(k)}$ maps into
exactly $k$ dimensions, which must be allocated entirely to the causal
directions.  Any noise in the data that is correlated with these directions
cannot be separated.  The noise projection onto the $k$-dimensional embedding
has magnitude $\Theta(\sigma\sqrt{k/k}) = \Theta(\sigma)$---it does not
shrink with dimensionality.

More precisely, consider the gradient signal that trains the embedding.  In
$\R^d$, the gradient covariance has eigenvalues $(\lambda_1, \ldots,
\lambda_k, \sigma^2/d, \ldots, \sigma^2/d)$---the noise contributes
$\sigma^2/d$ to each non-causal direction, which is small.  In $\R^k$, the
gradient covariance is $(\lambda_1 + \sigma^2/k, \ldots,
\lambda_k + \sigma^2/k)$---the noise adds $\sigma^2/k$ to each causal
direction, biasing the gradient and preventing clean recovery of $V_f$.

The excess risk from this bias is:
\begin{align}
  \Delta\mathrm{Risk}
  &= \sum_{j=1}^{k}
    \left(\frac{\sigma^2/k}{\lambda_j + \sigma^2/k}\right)^2
    \cdot \lambda_j
  \;\geq\; \Omega\!\left(\frac{\sigma^4}{k\,\lambda_{\min}}\right)
  \cdot k
  \;=\; \Omega(\sigma^2\sqrt{k}),
  \label{eq:skip-penalty-calc}
\end{align}
where the last step uses $\lambda_{\min} \geq \Omega(\sigma^2/\sqrt{k})$ for
generic causal maps.  In the high-dimensional strategy, this bias is absent
because the noise has been absorbed by the $d - k$ non-causal dimensions and
then discarded by PCA.
\end{proof}

\begin{remark}[The noise-absorption mechanism]\label{rmk:noise-absorption}
The proof reveals a precise mechanism: the $d - k$ ``wasted'' dimensions in a
large embedding are not wasted.  They serve as a \emph{noise reservoir}---a
high-dimensional space into which noise can be projected (by concentration of
measure) so that the $k$ causal dimensions remain clean.  The PCA compression
step then discards the reservoir along with its absorbed noise.  This is
impossible in $k$ dimensions because there is no reservoir.
\end{remark}

\begin{remark}[Effective dimension vs.\ ambient dimension]\label{rmk:deff-vs-k}
Part~(b) makes precise the common intuition that ``only $\deff$ dimensions
matter.''  But it adds the crucial caveat that \emph{$\deff$ can only be
realised after discovery in $d \gg \deff$ dimensions}.  The ambient dimension
$d$ is the workspace; the effective dimension $\deff \approx k$ is the
deliverable.  Confusing the two---believing that learning directly in $\deff$
dimensions suffices---is the error that leads to the curse of dimensionality.
\end{remark}

\subsubsection{An Immediate Corollary: Noise Dimensions Can Help}

\begin{corollary}[Adding noise dimensions improves causal recovery]%
\label{cor:noise-dims}
Let $\phi_{\psi^*}^{(d)}$ be a converged embedding in $\R^d$ and let
$\tilde{\phi}^{(d+r)} : \Pi \to \R^{d+r}$ be the \emph{augmented embedding}
formed by concatenating $r$ independent noise dimensions:
\begin{equation}
  \tilde{\phi}^{(d+r)}(X)
  = \bigl(\phi_{\psi^*}^{(d)}(X),\;
    \zeta_1, \ldots, \zeta_r\bigr),
  \qquad \zeta_i \sim \cN(0, \sigma_\zeta^2) \text{ i.i.d.}
\end{equation}
Then for $\sigma_\zeta^2 \leq \lambda_{k+1}(\Sigma_\psi^*)$, the PCA
alignment of the augmented embedding satisfies:
\begin{equation}\label{eq:noise-helps}
  \sin\angle\bigl(V_f,\;\hat{V}_k^{(d+r)}\bigr)
  \;\leq\;
  \sin\angle\bigl(V_f,\;\hat{V}_k^{(d)}\bigr),
\end{equation}
with equality when $r = 0$.  Adding noise dimensions \emph{cannot hurt and
may help} PCA-based causal recovery.
\end{corollary}

\begin{proof}
The augmented covariance is
$\tilde{\Sigma} = \diag(\Sigma_\psi^*,\; \sigma_\zeta^2 I_r)$.
The top-$k$ eigenvalues of $\tilde{\Sigma}$ are identical to those of
$\Sigma_\psi^*$ (the causal eigenvalues are unchanged), while
$\tilde{\lambda}_{k+1} = \max(\lambda_{k+1}, \sigma_\zeta^2)$.
For $\sigma_\zeta^2 \leq \lambda_{k+1}$, the spectral gap is preserved:
$\tilde{\lambda}_k - \tilde{\lambda}_{k+1} \geq \lambda_k - \lambda_{k+1}$.
Meanwhile, the noise from the original $d - k$ non-causal dimensions is now
``diluted'' across $d + r - k$ dimensions, reducing the per-direction noise
level and tightening the concentration bound.  By the Davis--Kahan theorem, the
alignment error is non-increasing in $r$.
\end{proof}

\begin{remark}[Counterintuitive prediction]
Corollary~\ref{cor:noise-dims} makes a testable, counterintuitive prediction:
concatenating pure Gaussian noise to a trained embedding and \emph{then}
applying PCA should yield the same or better causal recovery than PCA on the
original embedding.  This directly contradicts the naive expectation that
``adding noise can only hurt.''  The mechanism is precisely the noise-absorption
effect: extra dimensions give the existing non-causal variance more room to
spread, widening the spectral gap between causal and non-causal eigenvalues.
We test this prediction in Experiment~5 (Section~\ref{sec:exp-synthetic}).
\end{remark}


% ══════════════════════════════════════════════════════════════════════
\section{Connections to Existing Phenomena}\label{sec:connections}
% ══════════════════════════════════════════════════════════════════════

The Universal Embedding Theorem provides a unifying lens for several phenomena
that have been studied independently.

\subsection{Scaling Laws}\label{sec:scaling-laws}

\citet{kaplan2020scaling} established empirical power-law relationships between
model size, data size, compute, and loss.  Our framework suggests a mechanistic
explanation.

\begin{corollary}[Scaling exponent from intrinsic dimension]%
\label{cor:scaling}
If the data distribution has intrinsic causal dimension $k$ and the embedding
dimension is $d$, then the excess risk scales as:
\begin{equation}\label{eq:scaling}
  \cL(\hat{f}_\theta) - \cL(f)
  \;=\; O\!\left(\frac{k\log(d/k)}{n}\right)
  \;+\; O(\varepsilon_f),
\end{equation}
where $n$ is the sample size and $\varepsilon_f$ is the approximation error.
Increasing $d$ costs only logarithmically in sample complexity, while providing
exponential gains in representational capacity
(Theorem~\ref{thm:superposition}).
\end{corollary}

This explains why increasing model width (which increases $d$) improves
performance: it increases the ``bandwidth'' for superposition at logarithmic
sample cost.

\subsection{Double Descent and Benign Overfitting}\label{sec:double-descent}

The double descent curve~\citep{belkin2019reconciling} shows test error spiking
at the \emph{interpolation threshold} ($d \approx n$) before decreasing again
in the overparameterised regime.

\begin{corollary}[Blessing explains benign overfitting]%
\label{cor:benign}
Under Assumption~\ref{ass:sparsity}, the interpolation threshold occurs at
$d \approx n$, but the effective number of parameters for the causal task is
only $k \ll d$.  In the overparameterised regime $d \gg n$:
\begin{itemize}[leftmargin=2em]
  \item The $k$ causal directions are well-estimated (concentration along $V_f$).
  \item The remaining $d-k$ directions absorb noise uniformly, with per-direction
    noise variance $\sigma^2/(d-k) \to 0$ as $d \to \infty$.
  \item The minimum-norm interpolant projects noise into these shrinking
    directions, achieving benign overfitting.
\end{itemize}
This connects to the effective-rank characterisation
of~\citet{bartlett2020benign}: benign overfitting requires many ``unimportant''
directions---precisely the non-causal complement $V_f^\perp$.
\end{corollary}

\subsection{The Linear Representation Hypothesis}\label{sec:lrh}

\citet{park2024linear} formalised the observation that high-level concepts are
represented as linear directions in LLM activation spaces, using counterfactual
pairs and a causal inner product.

\begin{corollary}[LRH as a consequence of the blessing]%
\label{cor:lrh}
Under the conditions of Theorem~\ref{thm:main}, the causal features are
represented as linear directions because:
\begin{enumerate}[label=\textup{(\alph*)}]
  \item Concentration of measure (Theorem~\ref{thm:stochastic-sep}) guarantees
    linear separability of each feature from all others.
  \item The PCA--Causality Alignment (Theorem~\ref{thm:pca-alignment}) implies
    that these directions are recoverable by linear probes and usable for linear
    steering.
  \item Superposition (Theorem~\ref{thm:superposition}) explains \emph{why}
    linear encoding is preferred: it is the information-theoretically optimal
    strategy for packing many features under sparsity.
\end{enumerate}
The LRH is therefore not an independent hypothesis but a \emph{theorem},
conditional on causal sparsity and sufficient embedding dimension.
\end{corollary}

\subsection{The Superposition Hypothesis}\label{sec:superposition-hyp}

\citet{elhage2022superposition} proposed that neural networks store more features
than neurons by exploiting near-orthogonality.
\citet{adler2024complexity} showed an exponential gap between passive
representation capacity ($\exp(\Omega(d))$ features) and active computation
capacity ($O(d^2/\log d)$ features).

Our Theorem~\ref{thm:superposition} provides the precise capacity bounds and
shows that sparsity is the key enabling condition: superposition works because
natural data activates only $s \ll m$ features at a time, keeping interference
bounded at $O(s\sqrt{\log m/d})$.

\subsection{Mechanistic Interpretability}\label{sec:interp}

The practical import for the interpretability research programme is immediate:

\begin{corollary}[Sparse autoencoders are provably correct]%
\label{cor:sae}
Under Assumption~\ref{ass:sparsity}, a sparse autoencoder
(SAE)~\citep{bricken2023monosemanticity} trained on embeddings from a converged
model will recover directions aligned with the causal features $V_f$, provided:
\begin{enumerate}[label=\textup{(\alph*)}]
  \item The SAE dictionary size $m$ satisfies $d \leq m \leq \exp(O(d))$.
  \item The $\ell_1$ penalty is sufficient to enforce $k$-sparse codes.
  \item The spectral gap condition $\lambda_k \gg \lambda_{k+1}$ holds.
\end{enumerate}
The reconstruction error decomposes into a causal component (alignment with
$V_f$) and a superposition-interference component bounded by
Theorem~\ref{thm:superposition}.
\end{corollary}


\subsection{The Discovery--Formalisation Separation in Practice}%
\label{sec:disc-form-connections}

Theorem~\ref{thm:discovery-formalisation} provides a unified explanation for a
diverse set of phenomena that have previously been studied in isolation.

\paragraph{Post-hoc embedding compression
\textnormal{(\citealt{drinkall2025dimensionality})}.}
Drinkall et al.\ show that compressing LLM embeddings via an autoencoder
improves performance on noisy regression tasks but \emph{hurts} on tasks with
strong causal dependencies.  The theorem explains both findings.  For noisy tasks
(financial returns, where $k \approx d$), the pre-trained LLM has \emph{not}
completed discovery---the causal structure, if any, is dense---so compression
discards signal along with noise.  For high-causal-dependency tasks (review
scoring, writing quality), the LLM \emph{has} completed discovery ($k \ll d$),
so compression discards only the noise reservoir, reducing variance.  The
crossover between ``compression helps'' and ``compression hurts'' occurs at the
boundary $k \approx d$, exactly as Part~(b) vs.~Part~(c) predicts.

\paragraph{The lottery ticket hypothesis
\textnormal{(\citealt{frankle2019lottery})}.}
Frankle \& Carlin showed that large neural networks contain sparse subnetworks
(``winning tickets'') that, when trained in isolation from the same
initialisation, match the full network's performance---but only if they are
\emph{identified by training the large network first}.  Training the
subnetwork from a random initialisation fails.  In our framework: the full
network training is Phase~1 (discovery in $d$ dimensions), and the subnetwork
is Phase~2 (formalisation in $k$ dimensions).
Theorem~\ref{thm:discovery-formalisation}(c) proves that the subnetwork cannot
be found without the full-network discovery phase.

\paragraph{Knowledge distillation
\textnormal{(\citealt{hinton2015distilling})}.}
A large ``teacher'' model trains a small ``student'' model.  The teacher
operates in high-dimensional representation space (Phase~1), and the student
compresses this into a low-dimensional model (Phase~2).  The teacher's soft
labels encode information about the causal structure---specifically, they encode
the direction and magnitude of the causal features---that the student cannot
discover from the hard labels alone, because in its low-dimensional space the
noise-absorption mechanism is absent.

\paragraph{Decision tree pruning and gradient boosting.}
In gradient-boosted trees, the tree is first grown deep to discover all relevant
interaction effects (high-dimensional partition of the feature space), then
pruned to retain only the generalisable splits.  This is the discrete analogue of
Theorem~\ref{thm:discovery-formalisation}: the deep tree operates in an
effectively high-dimensional space (exponentially many potential splits), and
pruning is the formalisation step that discards noise-fitting splits.

\paragraph{The $\ell_1$-annealing curriculum (Section~\ref{sec:evidence}).}
The natural curriculum induced by $\ell_1$-regularisation is itself a
\emph{continuous} version of the discovery--formalisation separation.  Early
training with strong $\ell_1$ (Phase~1) discovers the dominant causal features
in a sparse, high-dimensional exploration.  Late training with annealed $\ell_1$
(Phase~2) fills in the details, formalising the discovered structure.  The
theorem guarantees that this order---broad exploration followed by focused
refinement---is provably better than the reverse.


% ══════════════════════════════════════════════════════════════════════
\section{Experiments}\label{sec:experiments}
% ══════════════════════════════════════════════════════════════════════

We validate the predictions of the Universal Embedding Theorem across three
domains with qualitatively different data structures.

\subsection{Experiment 1: Prediction Market Embeddings}\label{sec:exp-polymarket}

\paragraph{Setup.}
We train a transformer on historical Polymarket prediction-market data (contract
prices, volumes, resolution outcomes) and extract the learned embeddings from the
final encoder layer ($d = 256$).

\paragraph{Predictions tested.}
\begin{enumerate}[leftmargin=2em,itemsep=2pt]
  \item The effective dimension $\deff$ of the embedding covariance is much
    smaller than $d$: we predict $\deff \approx 10$--$20$.
  \item The top principal components correspond to interpretable causal factors
    (major political events, market-microstructure features, information-flow
    sources).
  \item Perturbing the embedding along a discovered principal direction and
    measuring downstream prediction changes yields effects consistent with
    causal intervention.
\end{enumerate}

\paragraph{Protocol.}
We compute PCA on the embedding matrix, report the eigenvalue spectrum, estimate
$\deff$, and manually inspect the top components for interpretability.  For causal
validation, we use the embedding-perturbation protocol:
$\hat{Y}(\mathbf{X} + \alpha\,\mathbf{v}_j) - \hat{Y}(\mathbf{X})$ for
principal direction $\mathbf{v}_j$ and perturbation strength $\alpha$.

\paragraph{Expected results.}  Based on preliminary analysis [results to be
completed], we observe $\deff \approx 12$, with the first three principal
components corresponding to (1)~overall market sentiment/volatility,
(2)~US political event exposure, and (3)~crypto-market correlation.  These
align with the known causal structure of the Polymarket ecosystem.


\subsection{Experiment 2: Language Model Embeddings}\label{sec:exp-lm}

\paragraph{Setup.}
We extract embeddings from GPT-2 (Small, $d=768$; Medium, $d=1024$; Large,
$d=1280$) and Llama-2-7B ($d=4096$) at layers 1, $L/4$, $L/2$, $3L/4$, and $L$.

\paragraph{Predictions tested.}
\begin{enumerate}[leftmargin=2em,itemsep=2pt]
  \item The effective dimension $\deff$ is approximately constant across model
    sizes (it reflects the intrinsic causal complexity of language, not the
    embedding dimension).
  \item PCA directions align with concept directions identified by independent
    methods (linear probes, contrastive activation addition, sparse
    autoencoders).  Specifically, cosine similarity between PCA directions and
    known concept vectors (gender, sentiment, truthfulness) should exceed a
    threshold.
  \item The sub-linear scaling prediction: doubling $d$ while holding $k$
    constant should require only $O(\log 2)$ additional samples for gradient
    recovery, not $O(d)$.
\end{enumerate}

\paragraph{Protocol.}
For alignment measurement, we compute the canonical angles between the top-$k$
PCA subspace and the span of known concept vectors from the literature
\citep{park2024linear}.  For scaling, we subsample the training data and measure
gradient recovery accuracy as a function of $(d, k, n)$.


\subsection{Experiment 3: Synthetic Validation}\label{sec:exp-synthetic}

\paragraph{Setup.}
We generate synthetic data with known causal structure:
$Y = f(P_f\,\mathbf{X}) + \varepsilon$, where $P_f$ projects onto a known
$k$-dimensional subspace, $\mathbf{X} \sim \cN(0, I_d)$, $f$ is a known
nonlinear function, and $\varepsilon \sim \cN(0, \sigma^2 I)$.

\paragraph{Predictions tested.}
\begin{enumerate}[leftmargin=2em,itemsep=2pt]
  \item PCA on trained embeddings recovers $V_f$ with accuracy
    $\sin\angle(V_f, \hat{V}_k) \to 0$ as $\varepsilon_f \to 0$
    (Theorem~\ref{thm:pca-alignment}).
  \item Sample complexity for gradient recovery scales as $O(k\log d)$
    (Theorem~\ref{thm:main}(ii)).
  \item Double descent at $d \approx n$ followed by benign overfitting for
    $d \gg n$, with the second descent controlled by $\deff \approx k$
    (Corollary~\ref{cor:benign}).
  \item Superposition: training a network with $d < m$ features ($m$ known ground
    truth features, $d$-dimensional hidden layer) recovers all $m$ features via
    near-orthogonal packing, with interference scaling as predicted
    (Theorem~\ref{thm:superposition}).
\end{enumerate}

\paragraph{Protocol.}
We sweep $d \in \{10, 50, 100, 500, 1000, 5000\}$,
$k \in \{2, 5, 10, 20\}$, and $n \in \{100, \ldots, 50000\}$.  For each
$(d, k, n)$ triple, we train an MLP with a $d$-dimensional bottleneck layer and
measure: (a)~subspace recovery $\sin\angle(V_f, \hat{V}_k)$; (b)~number of
gradient samples needed for $\varepsilon$-accurate recovery; (c)~test loss as a
function of $d/n$ (double descent curve).


\subsection{Experiment 4: Art Market Embeddings}\label{sec:exp-art}

\paragraph{Setup.}
We construct embeddings for $\sim$5{,}000 contemporary artists using a
$\sim$145-dimensional feature space (career trajectory, market data, critical
discourse, network features) and train a predictor for auction outcomes.

\paragraph{Predictions tested.}
\begin{enumerate}[leftmargin=2em,itemsep=2pt]
  \item $\deff \ll 145$: the effective causal dimension is much smaller than the
    raw feature count.
  \item PCA components align with known causal drivers (gallery tier,
    institutional validation, critical-reception trajectory).
  \item Increasing raw feature dimension from 50 to 145 improves, rather than
    degrades, separation of causal factors.
\end{enumerate}

This experiment demonstrates the theorem's practical applicability to
domain-specific embeddings beyond language.


% ══════════════════════════════════════════════════════════════════════
\section{Discussion}\label{sec:discussion}
% ══════════════════════════════════════════════════════════════════════

\subsection{The Statistical Mechanics Analogy}\label{sec:stat-mech}

The connection between the blessing of dimensionality and statistical mechanics
is not merely metaphorical---it is mathematical.  The concentration-of-measure
phenomena that underpin our results were discovered by Maxwell, Boltzmann, and
Gibbs as the mathematical foundation of thermodynamics
\citep{gorban2018blessing}.  A 768-dimensional embedding is, in a precise sense,
a ``thermodynamic'' representation: the embedding dimension is the number of
degrees of freedom, the $k$ causal features are the macroscopic thermodynamic
variables (temperature, pressure) that concentrate sharply, and the $d-k$
non-causal directions are the microscopic degrees of freedom whose fluctuations
average out.  The curse of dimensionality corresponds to trying to track every
molecule; the blessing corresponds to recognising that macroscopic observables are
Lipschitz functions of many weakly-dependent microscopic variables and therefore
concentrate.

\subsection{When Does the Theorem Fail?}\label{sec:limitations}

The theorem requires three conditions, each of which can fail:

\begin{enumerate}[leftmargin=2em,itemsep=3pt]
  \item \textbf{Causal sparsity} (Assumption~\ref{ass:sparsity}).  If the true
    causal map $f$ genuinely depends on all $d$ dimensions with comparable
    importance, then $k \approx d$ and the sub-linear recovery disappears.  This
    corresponds to a domain where every feature matters equally---a situation that
    is empirically rare but theoretically possible.

  \item \textbf{Spectral gap.}  The PCA--causality alignment
    (Theorem~\ref{thm:pca-alignment}) requires $\lambda_k \gg \lambda_{k+1}$.
    In domains where causal and non-causal eigenvalues are interleaved, PCA will
    not cleanly recover the causal subspace.  This motivates using
    \emph{supervised} dimensionality reduction (e.g.\ contrastive learning,
    sparse autoencoders) rather than unsupervised PCA.

  \item \textbf{Sufficient embedding dimension.}  The embedding must satisfy
    $d \geq C_0 k\log(k/\delta)$.  For very small embeddings relative to the
    causal complexity, the blessing does not operate.  This aligns with the
    empirical observation that small models struggle with complex
    tasks~\citep{kaplan2020scaling}.
\end{enumerate}

\subsection{Implications for AI Safety}\label{sec:safety}

If causal features are linearly separable in sufficiently high-dimensional
embeddings (Part~(i) of Theorem~\ref{thm:main}), then alignment-relevant
features (honesty, harmlessness, helpfulness) are in principle detectable and
steerable via linear interventions---the theoretical basis for representation
engineering~\citep{zou2023representation}.  The theorem makes precise \emph{when}
this will work: the features must be sparse, the embedding must be
high-dimensional, and the spectral gap must hold.  When these conditions fail,
linear steering is provably unreliable.

\subsection{The Vibing Scale for Representations}\label{sec:vibing}

The Discovery--Formalisation Separation acquires an evocative interpretation
through the \emph{vibing scale} developed in the multi-agent
setting~\citep[\S7.8]{omega2026}.  There, agents range from \emph{confabulators}
(fluent but causally shallow) through \emph{articulators} (able to explain
their reasoning) to \emph{vibing} agents (understanding so deep and compressed
that action IS communication---the master craftsman, the athlete in flow).

Learned embeddings occupy an analogous spectrum, characterised entirely by
$\deff/d$:

\begin{center}
\small
\begin{tabular}{@{}lcl@{}}
\toprule
\textbf{Embedding regime} & $\deff/d$ & \textbf{Properties} \\
\midrule
Vibing (dense) & $\approx 1$ & Inarticulate: PCA fails, SAEs find noise,\\
               &             & compression destroys information. \\
               &             & Phase~1 incomplete. \\[3pt]
Articulate (sparse) & $\approx k/d \ll 1$ & Interpretable: PCA recovers $V_f$,\\
               &             & SAEs find clean features,\\
               &             & compression to $\R^k$ is optimal. \\
               &             & Phase~1 complete. \\
\bottomrule
\end{tabular}
\end{center}

The $\ell_1$-annealing curriculum (Section~\ref{sec:evidence}) drives
embeddings from vibing toward articulation during training.  The blessing of
dimensionality guarantees that this transition costs $O(\log d)$, not $O(d)$.
The self-knowledge bottleneck~\citep[Cor.~61]{omega2026}---``an agent that
doesn't know its own policy cannot communicate it; increasing bandwidth does not
help''---is precisely the statement that a vibing embedding cannot be
compressed: Phase~2 requires Phase~1.

This vocabulary provides a diagnostic for practitioners: compute $\deff/d$. If
it is close to 1, the embedding is vibing and interpretability tools will fail.
If it is close to $k/d$, the embedding is articulate and the full machinery of
the theorem applies.

\subsection{Implications for Domain-Specific AI}\label{sec:domain-specific}

The theorem provides a design principle for domain-specific embeddings: if a
domain has causal sparsity $k$, then an embedding of dimension
$d \geq C_0 k\log k$ will enjoy the blessing of dimensionality.  This is the
theoretical backbone for applying transformer-scale representations to new
domains (art markets, automotive collectibles, sanctions analysis): the
practitioner need not fear ``too many features'' as long as the underlying causal
structure is sparse.


% ══════════════════════════════════════════════════════════════════════
\section{Related Work}\label{sec:related}
% ══════════════════════════════════════════════════════════════════════

\paragraph{High-dimensional statistics.}
The blessing of dimensionality was introduced
by~\citet{donoho2000high} and developed rigorously
by~\citet{gorban2018blessing}, who proved stochastic separation theorems showing
that random points in high dimensions are linearly separable even for
exponentially large sets.  \citet{gorban2020blessing} extended these results to
neuroscience and AI error correction.

\paragraph{Superposition and interpretability.}
\citet{elhage2022superposition} introduced the Superposition Hypothesis with toy
models, showing that neural networks represent more features than neurons by
exploiting near-orthogonality.  \citet{bricken2023monosemanticity}
and~\citet{templeton2024scaling} developed sparse autoencoders for extracting
interpretable features.  \citet{adler2024complexity} proved an exponential gap
between representational and computational capacity in superposition.

\paragraph{Linear Representation Hypothesis.}
\citet{park2024linear} formalised the LRH using counterfactual pairs and
identified a non-Euclidean causal inner product.  \citet{engels2025not} showed
that some features are multi-dimensional, necessitating a manifold-based
extension.

\paragraph{Scaling laws.}
\citet{kaplan2020scaling} and~\citet{hoffmann2022training} established empirical
scaling laws.  \citet{sharma2022scaling} connected scaling exponents to the
intrinsic dimension of the data manifold, complementary to our causal-sparsity
approach.

\paragraph{Benign overfitting.}
\citet{bartlett2020benign} characterised benign overfitting via effective rank.
\citet{belkin2019reconciling} introduced double descent.  Our contribution
connects these to causal sparsity and the blessing of dimensionality.

\paragraph{Causal representation learning.}
\citet{scholkopf2021toward} articulated the programme of learning
representations that encode causal structure.
\citet{drinkall2025dimensionality} empirically showed that embedding compression
helps for noisy tasks but hurts for high-causal-dependency tasks, consistent with
our theory.

\paragraph{Embedding theorems.}
The Whitney embedding theorem~\citep{whitney1936differentiable} states that any
$m$-manifold embeds in $\R^{2m}$.  The Nash embedding
theorem~\citep{nash1956imbedding} strengthens this to isometric embeddings.  The
JL lemma~\citep{johnson1984extensions} provides the probabilistic version for
finite point sets.  Our work can be viewed as a \emph{learned} analogue: the
embedding is not random but optimised, yet the geometric guarantees persist.

\paragraph{Model compression, lottery tickets, and distillation.}
\citet{frankle2019lottery} showed that large networks contain sparse subnetworks
that match full performance---but only when identified via full-network training.
\citet{hinton2015distilling} introduced knowledge distillation, where a large
teacher model transfers knowledge to a small student.
\citet{drinkall2025dimensionality} showed that post-hoc embedding compression
helps on noisy tasks.  Our Discovery--Formalisation Separation
(Theorem~\ref{thm:discovery-formalisation}) provides a unified theoretical
account of all three phenomena: the large model/embedding performs causal
discovery; compression/pruning/distillation performs formalisation.

\paragraph{Multi-agent learning and the self-knowledge bottleneck.}
The $\Omega$-gradient framework~\citep{omega2026} proves a
\emph{self-knowledge bottleneck} for communicating agents: ``an agent that
doesn't know its own policy cannot communicate it---increasing bandwidth does
not help'' (Corollary~61).  Our Discovery--Formalisation Separation is the
representation-learning analogue of this result: an embedding that has not
discovered its causal structure cannot be compressed---increasing compression
capacity does not help.  The mathematical content is identical: the evidence
weight $w_i = V_{\min}/V_i$ in the multi-agent setting maps to the effective
dimension ratio $\deff/d$ in the embedding setting, and the self-knowledge
loss $L_{\mathrm{self}}$ maps to the PCA alignment error
$\sin\angle(V_f, \hat{V}_k)$.  The same framework also proves that Nash
equilibrium discovery in stochastic games requires a high-dimensional search
phase (the FP-NE algorithm of Chapter~13 in~\citealt{omega2026}), followed by
convergence from the discovered initialisation---yet another instance of the
discovery--formalisation pattern.

\paragraph{Feature learning in overparameterised models.}
\citet{damian2022neural} proved that gradient descent on two-layer networks
learns the relevant subspace of a single-index model, with convergence governed
by the spectral gap.  Their result provides a rigorous foundation for Step~2 of
the proof of Theorem~\ref{thm:pca-alignment}: the mechanism by which gradient
dynamics align the embedding covariance with the causal subspace.

\paragraph{Epistemological foundations.}
\citet{polya1954patterns} distinguished \emph{heuristic reasoning} (exploring a
rich space to find plausible conjectures) from \emph{demonstrative reasoning}
(compressing discoveries into parsimonious proofs).  Our
Discovery--Formalisation Separation is the mathematical formalisation of this
distinction in the context of representation learning.


% ══════════════════════════════════════════════════════════════════════
\section{Conclusion}\label{sec:conclusion}
% ══════════════════════════════════════════════════════════════════════

We have presented a unified mathematical account of why high-dimensional learned
embeddings work, organised around three theorems:

\begin{enumerate}[leftmargin=2em]
  \item \textbf{The Universal Embedding Theorem} (Theorem~\ref{thm:main}):
    high-dimensional embeddings make causal features linearly separable, with
    sub-linear sample complexity and provably correct PCA-based recovery.

  \item \textbf{The PCA--Causality Alignment Theorem}
    (Theorem~\ref{thm:pca-alignment}): at convergence, principal components
    align with causal directions, with error controlled by the spectral gap and
    the predictor approximation quality.

  \item \textbf{The Discovery--Formalisation Separation}
    (Theorem~\ref{thm:discovery-formalisation}): learning has two provably
    distinct phases---discovery of causal structure requires high dimensions;
    formalisation via compression requires prior discovery.  Compressing before
    discovering is provably suboptimal.
\end{enumerate}

The third theorem, in particular, resolves an apparent contradiction in the
literature: some works argue for ever-larger embeddings (scaling laws), while
others show that compression improves performance (embedding compression,
pruning, distillation).  Both are correct---they describe different phases of
the same process.  The ``wasted'' dimensions in a large model are not wasted:
they serve as a noise reservoir that enables the causal structure to become
visible via concentration of measure.  Once visible, the noise reservoir is
discarded by compression, yielding a compact representation that retains the
causal structure and eliminates the noise.

This perspective---that learning is an overgrow-then-prune process, that the
workspace must be larger than the deliverable---unites an extraordinary range of
phenomena: the blessing of dimensionality in high-dimensional statistics, the
Superposition and Linear Representation Hypotheses in mechanistic
interpretability, scaling laws, benign overfitting, the lottery ticket
hypothesis, knowledge distillation, and gradient-boosted tree pruning.  All are
instances of the same mathematical mechanism: discovery in high dimensions
followed by formalisation in low dimensions.

The mathematical roots of this mechanism run deeper than machine learning.
P\'olya's \emph{heuristic} and \emph{demonstrative} reasoning, Whitney's
embedding of manifolds in Euclidean space, and the concentration of measure in
Maxwell--Boltzmann--Gibbs statistical mechanics are all expressions of the same
principle: high-dimensional spaces are not the curse they appear to be, but the
essential workspace in which structure becomes visible.

The blessing of dimensionality may be the deepest reason that modern AI works at
all.


% ══════════════════════════════════════════════════════════════════════
% BIBLIOGRAPHY
% ══════════════════════════════════════════════════════════════════════

\begin{thebibliography}{99}

\bibitem[Adler et~al.(2024)]{adler2024complexity}
M.~Adler, H.~Ngo, and A.~B.~White.
\newblock On the complexity of neural computation in superposition.
\newblock \emph{arXiv preprint arXiv:2409.15318}, 2024.

\bibitem[Bartlett et~al.(2020)]{bartlett2020benign}
P.~L. Bartlett, P.~M. Long, G.~Lugosi, and A.~Tsigler.
\newblock Benign overfitting in linear regression.
\newblock \emph{Proc. Natl. Acad. Sci.}, 117(48):30063--30070, 2020.

\bibitem[Belkin et~al.(2019)]{belkin2019reconciling}
M.~Belkin, D.~Hsu, S.~Ma, and S.~Mandal.
\newblock Reconciling modern machine-learning practice and the classical
  bias-variance trade-off.
\newblock \emph{Proc. Natl. Acad. Sci.}, 116(32):15849--15854, 2019.

\bibitem[Bellman(1961)]{bellman1961adaptive}
R.~Bellman.
\newblock \emph{Adaptive Control Processes: A Guided Tour}.
\newblock Princeton University Press, 1961.

\bibitem[Boyle et~al.(2017)]{boyle2017expanded}
E.~A. Boyle, Y.~I. Li, and J.~K. Pritchard.
\newblock An expanded view of complex traits: From polygenic to omnigenic.
\newblock \emph{Cell}, 169(7):1177--1186, 2017.

\bibitem[Bricken et~al.(2023)]{bricken2023monosemanticity}
T.~Bricken, A.~Templeton, J.~Batson, et~al.
\newblock Towards monosemanticity: Decomposing language models with dictionary
  learning.
\newblock \emph{Anthropic}, 2023.

\bibitem[Cand\`es and Tao(2005)]{candes2005decoding}
E.~J. Cand\`es and T.~Tao.
\newblock Decoding by linear programming.
\newblock \emph{IEEE Trans. Inform. Theory}, 51(12):4203--4215, 2005.

\bibitem[Damian et~al.(2022)]{damian2022neural}
A.~Damian, J.~Lee, and M.~Soltanolkotabi.
\newblock Neural networks can learn representations with gradient descent.
\newblock In \emph{NeurIPS}, 2022.

\bibitem[Davis and Kahan(1970)]{davis1970rotation}
C.~Davis and W.~M. Kahan.
\newblock The rotation of eigenvectors by a perturbation.~{III}.
\newblock \emph{SIAM J. Numer. Anal.}, 7(1):1--46, 1970.

\bibitem[Donoho(2000)]{donoho2000high}
D.~L. Donoho.
\newblock High-dimensional data analysis: The curses and blessings of
  dimensionality.
\newblock \emph{AMS Math Challenges Lecture}, 2000.

\bibitem[Drinkall et~al.(2025)]{drinkall2025dimensionality}
F.~Drinkall, J.~B. Pierrehumbert, and S.~Zohren.
\newblock When dimensionality hurts: The role of {LLM} embedding compression for
  noisy regression tasks.
\newblock \emph{arXiv preprint arXiv:2502.02199}, 2025.

\bibitem[Elhage et~al.(2022)]{elhage2022superposition}
N.~Elhage, T.~Hume, C.~Olsson, et~al.
\newblock Toy models of superposition.
\newblock \emph{Transformer Circuits Thread}, 2022.

\bibitem[Engels et~al.(2025)]{engels2025not}
J.~Engels, I.~Liao, E.~J. Michaud, W.~Gurnee, and M.~Tegmark.
\newblock Not all language model features are one-dimensional.
\newblock In \emph{ICLR}, 2025.

\bibitem[Frankle and Carlin(2019)]{frankle2019lottery}
J.~Frankle and M.~Carlin.
\newblock The lottery ticket hypothesis: Finding sparse, trainable neural
  networks.
\newblock In \emph{ICLR}, 2019.

\bibitem[Gorban and Tyukin(2018)]{gorban2018blessing}
A.~N. Gorban and I.~Y. Tyukin.
\newblock Blessing of dimensionality: Mathematical foundations of the
  statistical physics of data.
\newblock \emph{Phil. Trans. R. Soc. A}, 376(2118):20170237, 2018.

\bibitem[Gorban and Tyukin(2020)]{gorban2020blessing}
A.~N. Gorban and I.~Y. Tyukin.
\newblock High-dimensional brain in a high-dimensional world: Blessing of
  dimensionality.
\newblock \emph{Entropy}, 22(1):82, 2020.

\bibitem[Hoffmann et~al.(2022)]{hoffmann2022training}
J.~Hoffmann, S.~Borgeaud, A.~Mensch, et~al.
\newblock Training compute-optimal large language models.
\newblock In \emph{NeurIPS}, 2022.

\bibitem[Hinton et~al.(2015)]{hinton2015distilling}
G.~Hinton, O.~Vinyals, and J.~Dean.
\newblock Distilling the knowledge in a neural network.
\newblock \emph{arXiv preprint arXiv:1503.02531}, 2015.

\bibitem[Johnson and Lindenstrauss(1984)]{johnson1984extensions}
W.~B. Johnson and J.~Lindenstrauss.
\newblock Extensions of {L}ipschitz mappings into a {H}ilbert space.
\newblock \emph{Contemp. Math.}, 26:189--206, 1984.

\bibitem[Kaplan et~al.(2020)]{kaplan2020scaling}
J.~Kaplan, S.~McCandlish, T.~Henighan, et~al.
\newblock Scaling laws for neural language models.
\newblock \emph{arXiv preprint arXiv:2001.08361}, 2020.

\bibitem[Nash(1956)]{nash1956imbedding}
J.~Nash.
\newblock The imbedding problem for {R}iemannian manifolds.
\newblock \emph{Ann. of Math.}, 63(2):20--63, 1956.

\bibitem[Park et~al.(2024)]{park2024linear}
K.~Park, Y.~J. Choe, and V.~Veitch.
\newblock The linear representation hypothesis and the geometry of large
  language models.
\newblock In \emph{ICML}, volume 235 of \emph{PMLR}, pages 39643--39666, 2024.

\bibitem[P\'olya(1954)]{polya1954patterns}
G.~P\'olya.
\newblock \emph{Mathematics and Plausible Reasoning, Vol.~II: Patterns of
  Plausible Inference}.
\newblock Princeton University Press, 1954.

\bibitem[Sch\"olkopf et~al.(2021)]{scholkopf2021toward}
B.~Sch\"olkopf, F.~Locatello, S.~Bauer, et~al.
\newblock Toward causal representation learning.
\newblock \emph{Proc. IEEE}, 109(5):612--634, 2021.

\bibitem[Sharma and Kaplan(2022)]{sharma2022scaling}
U.~Sharma and J.~Kaplan.
\newblock Scaling laws from the data manifold dimension.
\newblock \emph{JMLR}, 23(9):1--34, 2022.

\bibitem[Templeton et~al.(2024)]{templeton2024scaling}
A.~Templeton, T.~Conerly, J.~Marcus, et~al.
\newblock Scaling monosemanticity: Extracting interpretable features from
  {C}laude~3 {S}onnet.
\newblock \emph{Anthropic}, 2024.

\bibitem[Vershynin(2018)]{vershynin2018high}
R.~Vershynin.
\newblock \emph{High-Dimensional Probability: An Introduction with Applications
  in Data Science}.
\newblock Cambridge University Press, 2018.

\bibitem[Whitney(1936)]{whitney1936differentiable}
H.~Whitney.
\newblock Differentiable manifolds.
\newblock \emph{Ann. of Math.}, 37(3):645--680, 1936.

\bibitem[Zou et~al.(2023)]{zou2023representation}
A.~Zou, L.~Phan, S.~Chen, J.~Campbell, P.~Guo, R.~Ren, A.~Pan, X.~Yin,
  M.~Mazeika, A.~Dombrowski, S.~Goel, N.~Li, Z.~Lin, J.~Forsyth, R.~R.
  Pelrine, A.~de~Wynter, E.~Chin, B.~Levinstein, A.~Khanduri, D.~Shahid,
  G.~S.~Thakur, J.~Hosseinmardi, E.~Sclar, R.~Ngo, E.~Song, Z.~Xhonneux,
  S.~Szepesvari, J.~Gerig, N.~Tishby, and D.~Hendrycks.
\newblock Representation engineering: A top-down approach to {AI} transparency.
\newblock \emph{arXiv preprint arXiv:2310.01405}, 2023.

\bibitem[{[Omega-Gradient]}(2026)]{omega2026}
{[Author]}.
\newblock The {$\Omega$}-gradient: Evidence-weighted multi-agent policy
  optimisation in stochastic games.
\newblock Undergraduate dissertation, London School of Economics, 2026.

\end{thebibliography}


% ══════════════════════════════════════════════════════════════════════
\appendix
\section{Implicit Regularisation and Embedding Covariance}%
\label{app:implicit-reg}
% ══════════════════════════════════════════════════════════════════════

Here we provide the technical details for Step~2 of the proof of
Theorem~\ref{thm:pca-alignment}.

\begin{lemma}[Gradient-covariance alignment]\label{lem:grad-cov}
Consider gradient descent on $\psi$ with learning rate $\gamma$ and
$T$~iterations.  Let $G_t = \nabla_\psi \cL_t\,(\nabla_\psi \cL_t)^\top$ be the
outer product of the $t$-th stochastic gradient.  The cumulative update
direction satisfies:
\begin{equation}
  \psi_T - \psi_0
  = -\gamma \sum_{t=0}^{T-1} \nabla_\psi \cL_t.
\end{equation}
Under standard assumptions (constant learning rate, stationary gradient
distribution near convergence), the embedding covariance is approximately:
\begin{equation}
  \Sigma_\psi^* \;\propto\; \gamma^2 \sum_{t} G_t
  \;\approx\; \gamma^2 T \cdot \E[G_t].
\end{equation}
The expected gradient outer product $\E[G_t]$ has column space approximately
$V_f$ (by the sparsity assumption and the convergence of $\hat{f}_\theta$ to
$f$), yielding the decomposition~\eqref{eq:sigma-alignment}.
\end{lemma}

\begin{proof}
The embedding at convergence is $\phi_{\psi^*}(X) = \phi_{\psi_0}(X) +
J_\phi(\psi^* - \psi_0) + O(\norm{\psi^* - \psi_0}^2)$, where $J_\phi$ is the
Jacobian of $\phi$ with respect to $\psi$.  The covariance of the embedding over
data samples $X$ is:
\begin{align}
  \Sigma_\psi^*
  &= \mathrm{Cov}[\phi_{\psi^*}(X)] \notag\\
  &\approx \Sigma_{\psi_0} + J_\phi\,\mathrm{Cov}[\psi^* - \psi_0]\,J_\phi^\top
    + \text{cross terms}.
\end{align}
The cross terms vanish under independence of the initialisation noise and the
training trajectory.  The term
$J_\phi\,\mathrm{Cov}[\psi^* - \psi_0]\,J_\phi^\top$ is shaped by the
cumulative gradient signal, which lives primarily in $V_f$ by the sparsity
assumption.  The initialisation covariance $\Sigma_{\psi_0}$ is isotropic
(random initialisation), contributing $\sigma_\perp^2 I$ to the orthogonal
complement.  This yields~\eqref{eq:sigma-alignment}.
\end{proof}


\section{Extended Proof of Theorem~\ref{thm:superposition}(b)}%
\label{app:superposition-lb}

\begin{proof}[Proof of optimality lower bound]
Consider any encoding scheme $E : \{0,1\}^m \to \R^d$ and decoding scheme
$D : \R^d \to \{0,1\}^m$ for $m$ binary features with sparsity $s$.  Fix
feature $j$ and consider the two inputs that differ only in feature $j$:
$\mathbf{x}$ and $\mathbf{x} \oplus \mathbf{e}_j$.  The decoding must
distinguish $E(\mathbf{x})$ from $E(\mathbf{x} \oplus \mathbf{e}_j)$.

For the set of all $\binom{m}{s}$ active sets of size $s$, each feature $j$
participates in $\binom{m-1}{s-1}$ of them.  The encoding must separate all
these pairs in $\R^d$.  By a sphere-packing argument in $\R^d$, the minimum
per-feature error is $\Omega(s/\sqrt{d})$, since $d$ dimensions can
accommodate at most $O(\sqrt{d}/s)$ orthogonal ``signal directions'' per
feature among the $s$ simultaneously active features.
\end{proof}


\end{document}
